%------------------------------------------------------------------------------%

\usepackage{calculo_varias_variaveis-1}

\title {Números Complexos}

%------------------------------------------------------------------------------%
\begin{document}

\addtitle

%------------------------------------------------------------------------------%
\section{Unidade Imaginária}

%------------------------------------------------------------------------------%
\begin{frame}{Motivação}
  Considere a equação
  \[ x^2 + 1 = 0 \]
  \pause
  Manipulando formalmente temos a solução
  \[ x = \sqrt{-1} \]
  \pause
  Porém, \(x^2 \geq 0\) para todo $x\in\R$

  \pause
  \vspace{\baselineskip}
  Não existe \(\sqrt{-1}\) dentro dos números reais
\end{frame}

%------------------------------------------------------------------------------%
\begin{frame}{Unidade Imaginária}
  \emph{Unidade imaginária}
  \[ i^2 = -1 \]

  \pause
  \vspace{\baselineskip}
  \(j^2=-1\) para a Engenharia Elétrica
\end{frame}

%------------------------------------------------------------------------------%
\begin{frame}{Potências de $i$}
  \begin{align*}
    \onslide<+->{i^1 & = i \\}
    \onslide<+->{i^2 & = -1 \\}
    \onslide<+->{i^3}
    \onslide<+->{& = i^2 \, i \;}
    \onslide<+->{ = (-1) i}
    \onslide<+->{ = -i}   \\
    \onslide<+->{i^4}
    \onslide<+->{ & = i^2 \, i^2}
    \onslide<+->{  = (-1)(-1)}
    \onslide<+->{ = 1}  \\
    \onslide<+->{i^5}
    \onslide<+->{ & = i^4 \, i \;}
    \onslide<+->{ = 1\cdot i}
    \onslide<+->{ = i} \\
    \onslide<+->{i^6}
    \onslide<+->{ & = i^4 \, i^2}
    \onslide<+->{ = 1(-1)}
    \onslide<+->{ = -1}    \\
    \onslide<+->{i^7}
    \onslide<+->{ & = i^4 \, i^3}
    \onslide<+->{ = 1(-i)}
    \onslide<+->{ = -i}    \\
    \onslide<+->{i^8}
    \onslide<+->{ & = i^4 \, i^4}
    \onslide<+->{  = 1\cdot 1}
    \onslide<+->{ = 1}
  \end{align*}
\end{frame}

%------------------------------------------------------------------------------%
\begin{frame}{Potências de $i$}
\begin{align*}
  \onslide<+->{i^0 & = 1\\}
  \onslide<+->{i^{-1}}
  \onslide<+->{& = \frac{1}{i}}
  \onslide<+->{  = \frac{i}{i^2}}
  \onslide<+->{  = \frac{i}{-1}}
  \onslide<+->{  = -i\\}
  \onslide<+->{i^{-2}}
  \onslide<+->{& = \frac{1}{i^2}}
  \onslide<+->{  = \frac{1}{-1}}
  \onslide<+->{  = -1\\}
  \onslide<+->{i^{-3}}
  \onslide<+->{& = \frac{1}{i^3}}
  \onslide<+->{  = \frac{1}{-i}}
  \onslide<+->{  = \frac{i}{-i^2}}
  \onslide<+->{  = \frac{i}{1}}
  \onslide<+->{  = i\\}
\end{align*}

\[
  \onslide<+->{i^{75} =}
  \onslide<+->{?}
\]
\end{frame}

%------------------------------------------------------------------------------%
\begin{frame}{Potências de $i$}
\begin{columns}
\column{0.5\linewidth}
  \onslide<+->{Para uma potência $n$ inteira}

  \onslide<+->{
  \vspace{\baselineskip}
  Dividimos $n$ por 4}
  \[
    \onslide<+->{
    \begin{array}{>{\;}c<{\;}>{\;}c<{\;}}
      n & \multicolumn{1}{|c}{4} \\ \cline{2-2}
      r & q
    \end{array}}
    \qquad
    \onslide<+->{
    \begin{array}{>{\;}c<{\;}>{\;}c<{\;}}
      75 & \multicolumn{1}{|c}{4} \\ \cline{2-2}
      3  & 18
    \end{array}}
  \]
  \onslide<+->{Assim}
  \begin{gather*}
    \onslide<+->{n = 4q + r} \\
    \onslide<+->{75 = 4\times 18 + 3}
  \end{gather*}
\column{0.5\linewidth}
  \begin{align*}
    \onslide<+->{i^n}
    \onslide<+->{& = i^{4q+r} \\}
    \onslide<+->{& = i^{4q} \; i^r \\}
    \onslide<+->{& = \left(i^4\right)^q \; i^r \\}
    \onslide<+->{& = \left(1\right)^q \; i^r \\}
    \onslide<+->{& = i^r}
  \end{align*}
  \onslide<+->{
  \[
    i^{75} = i^3 = -i
  \]}
\end{columns}
\end{frame}

%------------------------------------------------------------------------------%
\begin{frame}{Cuidado com Raízes}
\begin{columns}
\column{0.4\linewidth}
  Sabemos que
  \[
    \sqrt{a}\,\sqrt{b} = \sqrt{ab}
  \]

  \pause
  Por exemplo
  \[
    \sqrt{9}\,\sqrt{4}
    \pause
    = 3\times 2
    \pause
    = 6
  \]
  \[
    \pause
    \sqrt{9\times 4}
    \pause
    = \sqrt{36}
    \pause
    = 6
  \]
\column{0.6\linewidth}
  \pause
  \[
    \sqrt{-9}\,\sqrt{-4}
    \pause
    = 3i\times 2i
    \pause
    = 6i^2
    \pause
    = {\color<15->{red}-6}
  \]
  \[
    \pause
    \sqrt{(-9)(-4)}
    \pause
    = \sqrt{36}
    \pause
    = {\color<15->{red}6}
  \]

  \vspace{\baselineskip}
  \onslide<16>{\emph{Só é verdade para raízes reais!}}
\end{columns}
\end{frame}

%------------------------------------------------------------------------------%
%------------------------------------------------------------------------------%
\begin{question}
  Simplifique as expressões
  \[
    a = i^7
  \]

  \[
    b = i^{57}
  \]

  \[
    c = \frac{i^5 + i^{17}}{i^3}
  \]
\end{question}

%------------------------------------------------------------------------------%
\begin{resolution}{Solução}
  \[
    a
    = i^7
    \pause
    = i^4 \, i^3
    \pause
    = i^3
    \pause
    = i^2 i
    \pause
    = -i
  \]

  \pause
  \vspace{0.5\baselineskip}
  \[
    b
    = i^{57}
    \pause
    = i^{4\times 14} \, i^1
    \pause
    = i
  \]

  \pause
  \vspace{0.5\baselineskip}
  \[
    c
    = \frac{i^5 + i^{16}}{i^3}
    \pause
    = \frac{i^5}{i^3} + \frac{i^{16}}{i^3}
    \pause
    = i^{5-3} + i^{16-3}
    \pause
    = i^2 + i^{13}
    \pause
    = -1 + i^{4\times 3+1}
    \pause
    = -1 + i
  \]
\end{resolution}

%------------------------------------------------------------------------------%


%------------------------------------------------------------------------------%
\section{Números Complexos}

%------------------------------------------------------------------------------%
\begin{frame}{Números Complexos}

  Conjunto dos Números Complexos
  \[
    \mathbb{C} = \big\{\, z=a+bi \;\big|\; a, b \in \R \,\big\}
  \]

  \pause
  Parte real de $z$
  \[
    \Re(z) = a
  \]

  \pause
  Parte imaginária de $z$
  \[
    \Im(z) = b
  \]

  \pause
  Os números complexos com parte imaginária nula são os Reais
\end{frame}

%------------------------------------------------------------------------------%
\begin{frame}{Igualdade Entre Números Complexos}

  \(z_1=a+bi\) \; é \emph{igual} a \; \(z_2=c+di\)

  \pause
  \vspace{\baselineskip}
  se, e somente se,
  \[
     a = c
     \qquad\text{e}\qquad
     b = d
  \]
\end{frame}

%------------------------------------------------------------------------------%
\begin{frame}{Operações com os Números Complexos}
  \onslide<+->{Soma e subtração}
  \onslide<+->{
    \[
      (a+b\emph{i}) \;+\; (c+d\emph{i}) \;=\; (a+c) \;+\; (b+d) \emph{i}
    \]}
  \onslide<+->{
    \vspace{-\baselineskip}
    \[
      (a+b\emph{i}) \;-\; (c+d\emph{i}) \;=\; (a-c) \;+\; (b-d) \emph{i}
    \]
  }

  \vspace{-0.5\baselineskip}
  \onslide<+->{Produto}
  \begin{align*}
    \onslide<+->{
      (a+b\emph{i}) \cdot (c+d\emph{i})}
    \onslide<+->{
      & \;=\; a         \cdot (c+d\emph{i})
        \;+\; b\emph{i} \cdot (c+d\emph{i}) } \\
    \onslide<+->{
      & \;=\; ac + ad\emph{i} + bc\emph{i} + bd\emph{i}^2 } \\
    \onslide<+->{
      & \;=\; ac + ad\emph{i} + bc\emph{i} - bd } \\
    \onslide<+->{
      & \;=\; (ac-bd) \;+\; (ad+bc) \emph{i} }
  \end{align*}
\end{frame}

%------------------------------------------------------------------------------%
\begin{frame}{Conjugado}
  O \emph{conjugado} do número complexo
  \[
    z = a + bi
  \]
  \pause
  é o número
  \[
    \bar{z} = a - bi
  \]
\end{frame}

%------------------------------------------------------------------------------%
\begin{frame}{Propriedades dos Conjugados}
  \begin{enumerate}[<+->]
  \setlength\itemsep{1.5em}

  \item \;
  \(
    \bar{\bar{z}} \onslide<+->{= z}
  \)

  \item \; se \(\bar{z}=z\) \onslide<+->{então $z$ é real}

  \item \;
  \(
    z\bar{z}
    \onslide<+->{= (a+bi)(a-bi)}
    \onslide<+->{= a^2 -abi +abi - b^2 i^2}
    \onslide<+->{= a^2 + b^2}
  \)
  \end{enumerate}
\end{frame}

%------------------------------------------------------------------------------%
\begin{frame}{Operações com os Números Complexos}
  \onslide<+->{Divisão}%
  \onslide<+->{: um número $q$ é a divisão de $z_1=a+bi$ por $z_2=c+di$ se}
  \begin{columns}
  \column{0.5\baselineskip}
    \begin{align*}
      \onslide<+->{q z_2           & = z_1                   \\}
      \onslide<+->{q z_2 \bar{z}_2 & = z_1 \bar{z}_2         \\}
      \onslide<+->{q (c^2 + d^2)   & = (c+di)(a-bi)          \\}
      \onslide<+->{                & = ca -cbi +dai -dbi^2   \\}
      \onslide<+->{                & = ca -cbi +dai +dbi     \\}
      \onslide<+->{                & = (ac + bd) + (ad -bc)i}
    \end{align*}
  \column{0.5\baselineskip}
  \begin{align*}
    \onslide<+->{q & = \frac{(ac + bd) + (ad -bc)i}{c^2 + d^2}           \\[2mm]}
    \onslide<+->{  & = \frac{(ac + bd)}{c^2 + d^2} + \frac{(ad -bc)}{c^2 + d^2}i}
  \end{align*}
  \vspace{\baselineskip}
  \onslide<+->{
  \[
    \frac{z_1}{z_2} = \frac{(ac + bd)}{c^2 + d^2} + \frac{(ad -bc)}{c^2 + d^2}i
  \]}
  \end{columns}
\end{frame}

%------------------------------------------------------------------------------%
\begin{frame}{Corpo Complexo}
  A soma e o produto nos complexos tem as mesmas propriedades dos reais

  \vspace{\baselineskip}
  Por isso o conjunto dos complexos é chamado \emph{corpo}

  \vspace{\baselineskip}
  Exemplos de corpos: racionais, reais, complexos
\end{frame}

%------------------------------------------------------------------------------%
%------------------------------------------------------------------------------%
\begin{question}{Exemplo \theexemplo}
  Calcule
  \begin{gather*}
    a = \big[(3 + 2i) + (4 - 5i)\big](1 + i)                      \\[5mm]
    b = \big[ (2 + 3i) (1 - i) \big] \big[ (1 + 2i) (2 - i) \big] \\[2mm]
    c = \frac{3 + 4i}{2 - i}
  \end{gather*}
\end{question}

%------------------------------------------------------------------------------%
\begin{resolution}{Solução}
  \begin{align*}
  \onslide<+->{a & = \big[(3 + 2i) + (4 - 5i)\big](1 + i) \\}
  \onslide<+->{  & = [7 - 3i] (1 + i)                     \\}
  \onslide<+->{  & = 7 + 7i -3i -3i^2                     \\}
  \onslide<+->{  & = 7 + 7i -3i +3                        \\}
  \onslide<+->{  & = 10 + 4i                                }
  \end{align*}
\end{resolution}

%------------------------------------------------------------------------------%
\begin{resolution}{Solução}
  \begin{align*}
    \onslide<+->{b & = \big[ (2 + 3i) (1 - i) \big] \big[ (1 + 2i) (2 - i) \big] \\}
    \onslide<+->{  & = [ 2 - 2i + 3i + 3 ] [ 2 -i + 4i + 2 ]                     \\}
    \onslide<+->{  & = ( 5 + i ) ( 4 + 3i )                                      \\}
    \onslide<+->{  & = 20 + 15i + 4i -3                                          \\}
    \onslide<+->{  & = 17 + 19i                                                    }
  \end{align*}
\end{resolution}

%------------------------------------------------------------------------------%
\begin{resolution}{Solução}
  \begin{align*}
    \onslide<+->{c & = \frac{3 + 4i}{2 - i}                   \\}
    \onslide<+->{  & = \frac{(3 + 4i)(2 + i)}{(2 - i)(2 + i)} \\}
    \onslide<+->{  & = \frac{6 + 3i + 8i -4}{4 + 2i -2i +1}   \\}
    \onslide<+->{  & = \frac{2 + 11i}{5}                      \\}
    \onslide<+->{  & = \frac{2}{5} + \frac{11}{5}i              }
  \end{align*}
\end{resolution}

%------------------------------------------------------------------------------%


%------------------------------------------------------------------------------%
\section{Plano Complexo}

%------------------------------------------------------------------------------%
\begin{frame}<handout:0>[t]{Plano Complexo}
  \begin{gather*}
    z = x + yi \quad          \in \mathbb{C} \\
    (x, y)     \hspace{7.2ex} \in \R^2
  \end{gather*}
\end{frame}

%------------------------------------------------------------------------------%
\begin{frame}[t]{Plano Complexo}
  \begin{gather*}
    z = x + yi \quad          \in \mathbb{C} \\
    (x, y)     \hspace{7.2ex} \in \R^2
  \end{gather*}

  \addtext{9}{7,1.5}{\scalebox{0.7}[0.8]{\input{fig_complex_plane}}}

  \pause
  \begin{itemize}[<+->]
    \item Parte real       \tabto{32mm} \( \Re(z) = x \)     \\[4mm]
    \item Parte imaginária \tabto{32mm} \( \Im(z) = y \)     \\[4mm]
    \item Conjugado        \tabto{32mm} \( \bar{z} = x - yi \)        \\[4mm]
    \item Módulo           \tabto{32mm} \( \abs{z} = \rho = \sqrt{x^2+y^2\,}\) \\[4mm]
    \item Argumento        \tabto{32mm} \( \arg{(z)} = \varphi = \arctan\left(\frac{\,y\,}{x}\right)\)
  \end{itemize}
\end{frame}

%------------------------------------------------------------------------------%
\section{Polinômios}

%------------------------------------------------------------------------------%
\begin{frame}{Polinômios}
  Um \emph{polinômio} de grau $n\in\N$ de uma variável complexa $z$ \pause é
  \[
    P_n(z)
    = a_n z^n
    + a_{n-1} z^{n-1}
    + \cdots
    + a_2 z^2
    + a_1 z
    + a_0
  \]
  com \(z\), \(a_k\in\mathbb{C}\)

  \pause
  \vspace{\baselineskip}
  Um polinômio de grau $n$ possui $n$ raízes complexas

  \pause
  \vspace{\baselineskip}
  Se os coeficientes \(a_k\) forem reais, as raízes serão reais ou pares conjugados
\end{frame}
%------------------------------------------------------------------------------%
\begin{frame}{Raízes da Equação Quadrática}
  Raízes do polinômio
  \[
    ax^2 + bx + c = 0
    \qquad\qquad
    a\neq 0
    \qquad
    a, b, c \in \R
  \]

  \pause
  Fórmula de Bhaskara
  \[
    x = \dfrac{\displaystyle -b \pm \sqrt{\Delta}}{2a}
    \qquad\text{ onde }\quad
    \Delta = b^2 - 4 ac
  \]

  \pause
  Se \(\;\Delta<0\;\) temos
  \(\quad \sqrt{\Delta} = \sqrt{\abs{\Delta}}\;i \)
\end{frame}

%------------------------------------------------------------------------------%
%------------------------------------------------------------------------------%
\begin{question}
  Dado \(z=\frac{ \left(1-\sqrt{3}\right) + \left(1+\sqrt{3}\right)i}{2+2i} \)

  calcule

  a) a parte real de $z$

  b) a parte imaginária de $z$

  c) o módulo de $z$

  d) o argumento de $z$
\end{question}

%------------------------------------------------------------------------------%
\begin{resolution}{Avaliando $z$}
\begin{align*}
  z
  & = \frac{ \left(1-\sqrt{3}\right) + \left(1+\sqrt{3}\right)i}{2+2i} \\
  & = \frac{ \left[\left(1-\sqrt{3}\right) + \left(1+\sqrt{3}\right)i\right](2-2i)}{(2+2i)(2-2i)} \\
  & = \frac{ \left(1-\sqrt{3}\right)(2-2i) + \left(1+\sqrt{3}\right)(2i+2)}{2^2 + 2^2} \\
  & = \frac{ 2-2i-2\sqrt{3}+2\sqrt{3}i + 2i+2+2\sqrt{3}i +2\sqrt{3}}{8} \\
  & = \frac{ 4 + 4\sqrt{3}i}{8}
  \;=\; \frac{1}{2} + \frac{\sqrt{3}}{2}i
\end{align*}
\end{question}

%------------------------------------------------------------------------------%
\begin{resolution}{Partes real e imaginária de $z$}
  Parte real de $z$
  \[
    \Re(z) = \frac{1}{2}
  \]

  Parte imaginária de $z$
  \[
    \Im(z) = \frac{\sqrt{3}}{2}
  \]
\end{resolution}

%------------------------------------------------------------------------------%
\begin{resolution}{Módulo de $z$}
  \[
    \abs{z}
    = \sqrt{\left(\frac{1}{2}\right)^2 + \left(\frac{\sqrt{3}}{2}\right)^2}
    = \sqrt{\frac{1}{4} + \frac{3}{4}}
    = \sqrt{\frac{1+3}{4}}
    = \sqrt{1}
    = 1
  \]
\end{resolution}

%------------------------------------------------------------------------------%
\begin{resolution}{Argumento de $z$}
  \[
    \varphi=\arg(z)
  \]
  \[
    \sin(\varphi)
    = \frac{\Im(z)}{\abs{z}}
    = \frac{\sqrt{3}}{2}
  \]
  \[
    \cos(\varphi)
    = \frac{\Re(z)}{\abs{z}}
    = \frac{1}{2}
  \]
  \[
    \arg(z)
    = \ang{60}
    = \frac{\pi}{3} \text{rad}
  \]
\end{resolution}

%------------------------------------------------------------------------------%






%------------------------------------------------------------------------------%
\section{Lista Mínima}

%------------------------------------------------------------------------------%
\begin{frame}{Lista Mínima}
%%  Cálculo Vol. 2 do Thomas 12\textsuperscript{a} ed. -- Seção 14.8
%%  \begin{enumerate}
%%    \item Estudar o texto da seção
%%    \item Resolver os exercícios: 3, 5, 7, 12, 16, 18, 24 e 30
%%  \end{enumerate}

  \vfill
  Atenção: A prova é baseada no livro, não nas apresentações
\end{frame}

%------------------------------------------------------------------------------%
\end{document}
%------------------------------------------------------------------------------%
